%%%%%%%%%%%%%%%%%%%%%%%%%%%%%%%%%%%%%%%%%%%%%%%%%%%%%%%%%%%%%%%%%%%%%%%%%%%%%%
% root.tex — ReliaDINO RA-L submission (assembled 2026-07-15)
% Build: pdflatex root && bibtex root && pdflatex root && pdflatex root
%%%%%%%%%%%%%%%%%%%%%%%%%%%%%%%%%%%%%%%%%%%%%%%%%%%%%%%%%%%%%%%%%%%%%%%%%%%%%%
\documentclass[letterpaper, 10pt, conference]{ieeeconf}

\IEEEoverridecommandlockouts  % needed for \thanks
\overrideIEEEmargins          % needed to meet printer requirements

\usepackage{graphicx}
\usepackage{amsmath}
\usepackage{amssymb}
\usepackage{booktabs}
\usepackage{multirow}
\usepackage{xcolor}
\usepackage{tikz}
\usetikzlibrary{arrows.meta, positioning, calc, shapes.geometric}
\usepackage{url}
\usepackage{cite}

% ---- TODO marker (every unknown/missing item) ----
\newcommand{\todo}[1]{\textcolor{red}{[TODO: #1]}}

%%%%%%%%%%%%%%%%%%%%%%%%%%%%%%%%%%%%%%%%%%%%%%%%%%%%%%%%%%%%%%%%%%%%%%%%%%%%%%
% TITLE — option (1) chosen; alternatives kept below, commented.
%%%%%%%%%%%%%%%%%%%%%%%%%%%%%%%%%%%%%%%%%%%%%%%%%%%%%%%%%%%%%%%%%%%%%%%%%%%%%%
\title{\LARGE \bf
ReliaDINO: Frozen Vision Foundation Models with Per-Modality Low-Rank
Adaptation and Reliability-Anchored Fusion for Multimodal Semantic
Segmentation
}
% Alternative titles:
% \title{\LARGE \bf One Frozen Backbone, Four Sensors: Reliability-Anchored
%   Multimodal Semantic Segmentation with DINOv3 and Per-Modality LoRA}
% \title{\LARGE \bf Where Should Sensor Reliability Act? Calibrated Competence
%   Gating and Per-Class Routing for Multimodal Semantic Segmentation}

\author{\todo{Author names}% <-this % stops a space
\thanks{\todo{Funding / project support footnote.}}%
\thanks{\todo{$^{1}$Authors' affiliation, address, and e-mail (RA-L requires
full author block).}}%
}

\begin{document}

\maketitle
\thispagestyle{empty}
\pagestyle{empty}

%%%%%%%%%%%%%%%%%%%%%%%%%%%%%%%%%%%%%%%%%%%%%%%%%%%%%%%%%%%%%%%%%%%%%%%%%%%%%%
% ------------------------------------------------------------------
% 0_abstract.tex — body-only (\input into root.tex)
% ------------------------------------------------------------------
% TITLE OPTIONS (pick one in root.tex):
% (1) ReliaDINO: Frozen Vision Foundation Models with Per-Modality
%     Low-Rank Adaptation and Reliability-Anchored Fusion for
%     Multimodal Semantic Segmentation
% (2) One Frozen Backbone, Four Sensors: Reliability-Anchored
%     Multimodal Semantic Segmentation with DINOv3 and Per-Modality LoRA
% (3) Where Should Sensor Reliability Act? Calibrated Competence Gating
%     and Per-Class Routing for Multimodal Semantic Segmentation
% ------------------------------------------------------------------

\begin{abstract}
Semantic segmentation for robots operating in adverse conditions must fuse complementary sensors whose reliability varies with the environment, yet current state-of-the-art fusion methods depend on supervision beyond segmentation labels---ground-truth depth or condition annotations paired with text encoders---and on fusion backbones trained from scratch. We present \emph{ReliaDINO}, which instead multiplexes a single \emph{frozen} DINOv3 ViT-L/16 across RGB, depth, event, and LiDAR inputs through independent per-modality low-rank adapters ($\approx$3.1M parameters), exchanges information with two cross-modal attention layers, and applies self-derived reliability where class decisions are formed: a calibrated competence gate over modality features and a per-class reliability-anchored router over per-modality posteriors. Trained with segmentation labels only, ReliaDINO reaches \todo{final: currently 57.14, training in progress} mIoU on the DELIVER test benchmark, surpassing the strongest published four-modality method, DGFusion (56.71), which requires ground-truth log-depth supervision and condition meta-text. We further report a systematic negative result: across five model generations, injecting the same reliability signals as pre-softmax attention-logit biases is ineffective or harmful once per-modality adaptation aligns the representations---locating the output gate and router, not the attention logits, as the effective site for reliability. We also contribute the first per-class audit of the DELIVER test split and a strictly validation-based checkpoint-selection protocol.
\end{abstract}

% \begin{IEEEkeywords}
% Sensor Fusion, Semantic Scene Understanding, Deep Learning for Visual Perception, Computer Vision for Transportation
% \end{IEEEkeywords}


% ---------------- Fig. 1: teaser (supervision budget vs. test mIoU) --------
\begin{figure}[t]
  \centering
  \resizebox{0.98\linewidth}{!}{% ============================================================
% fig:teaser — supervision budget vs. DELIVER test mIoU (CLDE)
% BODY-ONLY: \input this file inside a single-column {figure}.
% Preamble needs: \usepackage{tikz, xcolor}
%   \usetikzlibrary{positioning, arrows.meta, shapes.geometric}
% Data: notes/fact_experiments.md (a) test-CLDE cluster + (0).
% y mapping: y[cm] = (mIoU - 52) * 0.9
% Suggested caption (place in the figure env):
%   \caption{Extra supervision vs.\ DELIVER test accuracy (mIoU, all four
%   modalities). Prior condition- and depth-guided fusion methods buy
%   robustness with additional supervision: DGFusion uses ground-truth
%   log-depth and condition meta-text, CAFuser uses condition annotations
%   with CLIP. ReliaDINO is trained with segmentation labels only and
%   surpasses the previous best test result (56.71). The open circle is our
%   no-router ablation variant under test-best checkpoint selection; see
%   Sec.~\todo{ref} for the val-only selection protocol.
%   \todo{Update star and margin once the final run completes.}}
% ============================================================
\definecolor{OursBlue}{RGB}{0,90,181}
\definecolor{BaseGray}{RGB}{115,115,115}
\definecolor{SupRed}{RGB}{165,50,45}
\begin{tikzpicture}[font=\scriptsize, line cap=round,
  >={Stealth[length=1.6mm]},
  basedot/.style={circle, fill=BaseGray, inner sep=1.5pt},
  supdot/.style={circle, fill=SupRed, inner sep=1.5pt},
  supdotl/.style={circle, fill=SupRed!60, inner sep=1.5pt},
  oursopen/.style={circle, draw=OursBlue, thick, fill=white, inner sep=1.5pt},
  oursstar/.style={star, star points=5, star point ratio=2.2,
                   fill=OursBlue, draw=OursBlue, inner sep=1.7pt}]

  % ---- gridlines + y axis ----
  \foreach \v/\y in {53/0.9, 54/1.8, 55/2.7, 56/3.6, 57/4.5, 58/5.4}{
    \draw[BaseGray!25, line width=0.3pt] (0.05,\y) -- (7.5,\y);
    \draw (0,\y) -- (-0.07,\y) node[left, inner sep=1pt] {\v};
  }
  \draw[->] (0,0) -- (0,5.85) node[above, inner sep=1pt] {\shortstack{test mIoU\\(CLDE)}};
  \draw (0,0) -- (7.6,0);

  % ---- prior test-SOTA line ----
  \draw[dashed, SupRed, line width=0.6pt] (0.05,4.24) -- (7.5,4.24);
  \node[SupRed, fill=white, inner sep=1pt, anchor=west] at (0.15,4.44)
    {prior test SOTA 56.7};

  % ---- column 1: GT depth + condition text ----
  \node[supdot] (dg) at (1.0,4.24) {};
  \node[anchor=north, inner sep=2.5pt] at (dg.south) {DGFusion};
  \node[align=center, text=SupRed] at (1.3,-0.55)
    {GT depth +\\condition text};

  % ---- column 2: condition labels + CLIP ----
  \node[supdotl] (cafa) at (2.9,2.88) {};
  \node[anchor=west, inner sep=2pt] at (cafa.east) {CAFuser-CAA 55.2};
  \node[supdotl] (caf) at (2.9,3.24) {};
  \node[anchor=west, inner sep=2pt] at (caf.east) {CAFuser 55.6};
  \node[align=center, text=SupRed!70] at (3.4,-0.55)
    {condition labels\\+ CLIP};

  % ---- column 3: segmentation labels only ----
  \node[basedot] (cmx) at (6.4,0.90) {};
  \node[anchor=east, inner sep=2pt] at (cmx.west) {CMNeXt 53.0};
  \node[basedot] (stf) at (6.4,1.26) {};
  \node[anchor=east, inner sep=2pt] at (stf.west) {StitchFusion 53.4};
  \node[basedot] (gmf) at (6.4,2.25) {};
  \node[anchor=east, inner sep=2pt] at (gmf.west) {GeminiFusion 54.5};
  \node[oursopen] (ournr) at (6.4,5.04) {};
  \node[anchor=east, inner sep=2pt, text=OursBlue] at (ournr.west)
    {ReliaDINO (no router) 57.6};
  \node[oursstar] (ours) at (6.4,4.63) {};
  \node[anchor=east, inner sep=2pt, text=OursBlue] at (ours.west)
    {\textbf{ReliaDINO} 57.1\,\todo{final}};
  \node[align=center, text=OursBlue] at (6.2,-0.55)
    {segmentation\\labels only};

  % ---- margin arrow ----
  \draw[->, OursBlue, thick] (7.15,4.24) -- (7.15,4.63)
    node[midway, right, inner sep=1.5pt] {+0.43};

  % ---- supervision direction ----
  \draw[->, BaseGray] (0.6,-1.15) -- (7.0,-1.15)
    node[midway, below, inner sep=1.5pt, text=BaseGray]
    {decreasing extra supervision};
\end{tikzpicture}
}
  \caption{Extra supervision vs.\ DELIVER test accuracy (mIoU, all four
  modalities). Prior condition- and depth-guided fusion methods buy
  robustness with additional supervision: DGFusion uses ground-truth
  log-depth and condition meta-text, CAFuser uses condition annotations
  with CLIP. ReliaDINO is trained with segmentation labels only and
  surpasses the previous best test result (56.71). The open circle is our
  no-router ablation variant under test-best checkpoint selection; see
  Sec.~\ref{sec:setup} for the val-only selection protocol.
  \todo{Update star and margin once the final run completes.}}
  \label{fig:teaser}
\end{figure}

% ------------------------------------------------------------------
% 1_intro.tex — body-only (\input into root.tex)
% ------------------------------------------------------------------
\section{Introduction}
\label{sec:intro}

Robots and autonomous vehicles must perceive their surroundings under night, fog, rain, and sensor failure---conditions in which no single sensor is dependable. Multimodal suites that pair cameras with depth, LiDAR, and event sensors offer complementary evidence, and benchmarks such as DELIVER~\cite{cmnext2023} and MUSES~\cite{muses2024} have made adverse-condition multimodal segmentation a measurable problem. The central difficulty is not only \emph{how} to fuse heterogeneous inputs but \emph{how much to trust each sensor at each location}: sensor competence varies per condition, per region, and---critically---per semantic class, so a fusion policy that is fixed or merely scene-global discards exactly the complementarity that motivates the sensor suite.

Current state-of-the-art systems answer this question with additional supervision and bespoke fusion backbones. DGFusion~\cite{dgfusion2026}, the strongest published method on the DELIVER test split, derives its reliability signal from ground-truth log-depth supervision combined with condition meta-text through a language--vision encoder; CAFuser~\cite{cafuser2025} requires condition annotations and CLIP-based~\cite{clip2021} contrastive training; HyperDUM~\cite{hyperdum2025} builds label-derived prototypes. Orthogonally, the dominant architecture line---CMX~\cite{cmx2023}, CMNeXt~\cite{cmnext2023}, StitchFusion~\cite{stitchfusion2024}, GeminiFusion~\cite{geminifusion2024}---trains full multimodal encoders end-to-end, forgoing the robustness of large frozen vision foundation models (VFMs) whose self-supervised features transfer strongly to dense prediction~\cite{dinov3_2025}.

We take the opposite route (Fig.~\ref{fig:teaser}). \emph{ReliaDINO} encodes all four modalities with a \emph{single frozen} DINOv3 ViT-L/16~\cite{dinov3_2025}; modality identity enters only through independent per-modality low-rank adapters~\cite{lora2022} on the attention projections ($\approx$3.1M parameters in total). Light per-modality decoders yield calibrated class posteriors from which all reliability signals are computed---without any supervision beyond the segmentation labels. Two shared cross-modal attention layers exchange information across modalities, and reliability is applied where class decisions are formed: a calibrated competence gate with a training-free veto floor fuses the modality features, and a per-class reliability-anchored router (PRR) re-weights the per-modality posteriors, letting each class rely on the sensor that actually sees it.

Perhaps our most instructive finding is negative. The seemingly natural injection site for reliability---an additive, pre-softmax bias on cross-modal attention logits, a placement adjacent to recent attention-modulation work~\cite{sam2long2024,sae2026,primed2026}---is consistently ineffective in our setting: across five model generations spanning SAM2-based~\cite{sam2_2024,memorysam2025} and DINOv3-based architectures, such biases changed predictions negligibly or hurt significantly, whereas the identical signals contribute measurably at the output gate and router. Our analysis attributes this to per-modality low-rank adaptation itself: once adapters align the modality representations, attention no longer routes pathologically, and suppressing unreliable keys becomes a no-op---raising \emph{where to inject reliability} as a design question in its own right.

Our contributions are:
\begin{itemize}
  \item \textbf{C1 --- Frozen-VFM modality multiplexing.} A single frozen DINOv3 ViT-L/16 serves all four modalities via per-modality LoRA; to our knowledge the first DINOv3-based multimodal segmentation model, with 46.8M trainable parameters (13.4\%)---fewer than the total of a CMNeXt-B2 baseline.
  \item \textbf{C2 --- Reliability at the decision level.} Calibrated competence-gated fusion (+0.26 test mIoU at full resolution) and a per-class reliability-anchored router; reliability acts at the output gate and router rather than as an attention-logit bias.
  \item \textbf{C3 --- Supervision minimality.} Segmentation labels only; no ground-truth depth, condition annotations, or text encoders, in contrast to the methods we surpass.
  \item \textbf{C4 --- Analysis.} A five-generation controlled negative result on attention-logit reliability biasing; the first per-class audit of the DELIVER \emph{test} split, exposing a benchmark ceiling (Wall, Bridge, and Other are near-zero even for the official baseline \todo{confirm official CMNeXt per-class test numbers, M0-a}); and a strictly validation-based checkpoint-selection protocol.
\end{itemize}

On DELIVER with all four modalities, ReliaDINO attains \todo{final test mIoU; best so far 57.14} on the test split (previous best: DGFusion, 56.71) and \todo{final val mIoU; best so far 67.74} on the validation split, while its no-router variant reaches 68.19 val / 57.60 test mIoU---evidence that a frozen foundation model, minimally adapted and fused by self-derived reliability, outperforms specialized, extra-supervised fusion architectures.

% ============================================================
% Section 2 — Related Work (body-only; \input into root.tex)
% Facts: notes/fact_related.md (taxonomy plan ii, contrasts, per-competitor table iii)
% All \cite keys exist in latex/references.bib.
% ============================================================

\section{Related Work}
\label{sec:related}

\subsection{Multimodal Semantic Segmentation}
\label{sec:related:mmss}

Benchmarks with synchronized heterogeneous sensors have driven this field: DELIVER~\cite{cmnext2023} pairs RGB with depth, event, and LiDAR data under five adverse conditions and corruption cases, MUSES~\cite{muses2024} provides multi-sensor driving scenes under visual degradation, and MCubeS~\cite{mcubes2022} targets material segmentation. On the architecture side, CMX~\cite{cmx2023} introduced cross-modal feature rectification and fusion for RGB-X segmentation, TokenFusion~\cite{tokenfusion2022} substitutes uninformative tokens across modalities inside a ViT, GeminiFusion~\cite{geminifusion2024} performs pixel-wise intra- and inter-modal attention, and StitchFusion~\cite{stitchfusion2024} weaves features between frozen pretrained encoders through lightweight adapters. A second line pursues \emph{arbitrary-modal} robustness: MAGIC and MAGIC++~\cite{magic2024,magicpp2024} counter RGB-centric bias with hierarchical modality selection, AnySeg~\cite{anyseg2024} distills uni- and cross-modal knowledge to survive missing sensors, and OmniSegmentor~\cite{omnisegmentor2025} pretrains a single encoder on a five-modality ImageNeXt corpus, reaching 68.0 mIoU on the DELIVER \emph{val} split---a pretraining axis that is orthogonal to, and composable with, our adaptation-based approach. Common to all of these methods is that they adapt \emph{representations or architectures} to multiple sensors; none of them decides at inference time \emph{how much each sensor should be trusted}.

\subsection{Condition- and Reliability-Aware Fusion}
\label{sec:related:reliability}

A more recent family conditions the fusion process on scene state or sensor quality. CAFuser~\cite{cafuser2025} derives a condition token from the RGB stream in CLIP~\cite{clip2021} space and uses it to steer fusion, which requires verbo-visual contrastive supervision and condition annotations. DGFusion~\cite{dgfusion2026}, the current DELIVER \emph{test} state of the art (56.7 mIoU, four-modality CLDE setting), treats LiDAR both as an input and as a supervision target: ground-truth depth trained with a robust log-depth loss yields local depth tokens and a global condition token that condition cross-modal fusion. HyperDUM~\cite{hyperdum2025} estimates deterministic uncertainty from hyperdimensional prototype distances---prototypes built from labels plus a fine-tuned weighting layer---and reweights features before fusion (66.30$\rightarrow$67.59 mIoU on DELIVER \emph{val}). Evidential approaches learn uncertainty heads: TMC~\cite{tmc2021} aggregates Dirichlet opinions for multi-view classification and UTFNet~\cite{utfnet2023} applies uncertainty-guided feature weighting to RGB-thermal segmentation. Reliability has also been injected at the loss or adaptation level: functional-entropy regularization balances modality contributions during training only~\cite{zheng2025entropy}, READ~\cite{read2024} reweights a test-time-adaptation loss against modality reliability bias in classification, and ReliFusion~\cite{relifusion2025} scales the \emph{output} of cross-attention by a learned reliability score for 3D detection. Across this family, the reliability or condition signal invariably (i) depends on supervision beyond segmentation labels---ground-truth depth, condition text or annotations, label-built prototypes, or dedicated evidential losses---and/or (ii) enters at the feature, attention-output, or loss level rather than where per-class decisions are formed.

\subsection{Foundation Models and Parameter-Efficient Adaptation}
\label{sec:related:vfm}

Parameter-efficient fine-tuning makes frozen vision foundation models practical for dense prediction: LoRA~\cite{lora2022} adds low-rank updates to frozen projections, ViT-Adapter~\cite{vitadapter2023} equips plain ViTs with dense-prediction inductive biases, and SAMed~\cite{samed2023} established the LoRA-on-SAM precedent. Building on SAM and SAM~2~\cite{sam2023,sam2_2024}, MemorySAM~\cite{memorysam2025} treats modalities as video ``frames'' fused by SAM~2 memory attention---the architectural ancestor of the attention-level reliability biasing that we audit and ultimately reject in Sec.~\ref{sec:experiments}. MLE-SAM~\cite{mlesam} attaches modality-specific mixture-of-LoRA experts with a learned dispatch gate to a frozen SAM encoder, and MM~SAM-adapter~\cite{mmsamadapter2025} injects fused features into a SAM ViT-L via deformable cross-attention, but supports only two modalities per model.\footnote{MM~SAM-adapter reports DELIVER test results under its own two-modality protocol (e.g., 57.35 mIoU RGB-D), which is not comparable to the four-modality CLDE setting used in Table~\ref{tab:sota}.} In contrast to the SAM lineage, DINOv3~\cite{dinov3_2025} provides a self-supervised backbone with strong frozen dense features; to our knowledge, as of mid-2026, no prior multimodal segmentation method builds on DINOv3.\todo{re-run the DINOv3-multimodal-seg arXiv sweep at submission time; nearest prior work uses DINOv2 with a click-based NoC metric rather than mIoU} These adaptation methods answer \emph{how to represent} each modality on a frozen backbone, but not \emph{how much to trust} it at a given pixel and class.

\medskip
\noindent\textbf{Positioning.}
ReliaDINO combines both threads under a minimal supervision budget: a single frozen DINOv3 ViT-L multiplexed by per-modality LoRA adapters, with reliability applied \emph{only} where class decisions are formed---a calibrated competence gate at the fusion output and a per-class reliability-anchored router---rather than as an attention-output scale~\cite{relifusion2025}, a TTA loss weight~\cite{read2024}, or a feature reweighting stage~\cite{hyperdum2025}. Unlike DGFusion~\cite{dgfusion2026} and CAFuser~\cite{cafuser2025}, it needs no ground-truth depth supervision, no condition meta-text, and no condition annotations: competence is self-derived from the model's own calibrated per-modality predictions, using segmentation labels only. Finally, we report a systematic negative result on the remaining placement option, \emph{reliability as a pre-softmax attention-logit bias}. Nearby instantiations of that cell exist---a learned modality-prior additive bias for referring audio-visual segmentation~\cite{primed2026}, training-free entropy-derived additive attention modulation in LVLM decoding~\cite{sae2026}, and training-free multiplicative key scaling in SAM~2 memory attention for long videos~\cite{sam2long2024}---but none addresses multimodal dense segmentation. Our five-generation audit (Sec.~\ref{sec:experiments}) indicates why: once per-modality LoRA on a strong frozen backbone absorbs modality adaptation, logit-level suppression becomes a no-op, and reliability information is useful only at the gate and router.


% ---------------- Fig. 2: architecture overview (two-column) ---------------
\begin{figure*}[t]
  \centering
  \resizebox{\textwidth}{!}{% ============================================================
% fig:arch — ReliaDINO full architecture overview
% BODY-ONLY: \input inside a two-column {figure*} environment.
% Preamble needs: \usepackage{tikz, xcolor, amsmath}
%   \usetikzlibrary{positioning, arrows.meta, calc, shapes.geometric}
% Recommended wrapping: \resizebox{\textwidth}{!}{% ============================================================
% fig:arch — ReliaDINO full architecture overview
% BODY-ONLY: \input inside a two-column {figure*} environment.
% Preamble needs: \usepackage{tikz, xcolor, amsmath}
%   \usetikzlibrary{positioning, arrows.meta, calc, shapes.geometric}
% Recommended wrapping: \resizebox{\textwidth}{!}{% ============================================================
% fig:arch — ReliaDINO full architecture overview
% BODY-ONLY: \input inside a two-column {figure*} environment.
% Preamble needs: \usepackage{tikz, xcolor, amsmath}
%   \usetikzlibrary{positioning, arrows.meta, calc, shapes.geometric}
% Recommended wrapping: \resizebox{\textwidth}{!}{\input{figures/fig_arch}}
% Structure source: notes/fact_method.md §0–§9 (no numbers to update).
% Suggested caption (place in the figure* env):
%   \caption{ReliaDINO. All four modalities are encoded by a single frozen
%   DINOv3 ViT-L/16; modality identity enters only through independent
%   low-rank adapters (color-coded, rank~8 on Q/V of all 24 blocks).
%   Per-modality auxiliary decoders provide the class posteriors from which
%   calibrated reliability $r_m$ and the veto signal $v_m$ are computed
%   (purple; detached). Modalities exchange information through two shared
%   cross-modal attention layers; the fused map is produced by a calibrated
%   competence gate with a training-free veto floor, decoded by a SimpleFPN
%   head. The per-class reliability-anchored router adds a residual routed
%   prediction scaled by a zero-initialized $\alpha$. Reliability acts only
%   at the gate and router: the attention-logit bias site is removed
%   (struck out; Sec.~\todo{ref}). Dashed gray parts are training-only.}
% ============================================================
\definecolor{cRGB}{RGB}{213,94,0}
\definecolor{cDep}{RGB}{0,114,178}
\definecolor{cEvt}{RGB}{230,159,0}
\definecolor{cLid}{RGB}{0,158,115}
\definecolor{cRel}{RGB}{136,78,160}
\definecolor{cGray}{RGB}{120,120,120}
\definecolor{OursBlue}{RGB}{0,90,181}
\begin{tikzpicture}[font=\scriptsize, >={Stealth[length=1.5mm]}, line cap=round,
  box/.style={draw, rounded corners=2pt, align=center, inner sep=3pt,
              minimum height=9mm},
  mod/.style={draw, very thick, minimum width=7.5mm, minimum height=6mm,
              inner sep=1pt, align=center, font=\tiny},
  lorachip/.style={draw, very thick, rounded corners=1.5pt, minimum width=14mm,
              minimum height=3.6mm, inner sep=1pt, font=\tiny, fill=white},
  card/.style={draw, fill=white, minimum width=6.5mm, minimum height=5mm,
              inner sep=0pt},
  loss/.style={draw, dashed, rounded corners=2pt, text=cGray, draw=cGray,
              align=center, inner sep=2.5pt, font=\tiny},
  sig/.style={->, cRel, thick},
  flow/.style={->, thick},
  opn/.style={circle, draw, thick, inner sep=0.5pt, font=\footnotesize}]

  % ================= inputs =================
  \node[mod, draw=cRGB] (in1) at (0,3.9)  {RGB};
  \node[mod, draw=cDep] (in2) at (0,3.0)  {Depth};
  \node[mod, draw=cEvt] (in3) at (0,2.1)  {Event};
  \node[mod, draw=cLid] (in4) at (0,1.2)  {LiDAR};
  \node[font=\tiny, text=cGray, align=center] at (0,0.55) {$x_m$: $4{\times}(3,1024^2)$};

  % ================= shared frozen encoder =================
  \node[box, minimum width=32mm, minimum height=34mm, anchor=west]
    (enc) at (1.1,2.55) {};
  \node[align=center, anchor=north] at (2.7,4.15)
    {\textbf{Frozen DINOv3 ViT-L/16}\\shared, 24 blocks};
  \node[font=\tiny, text=cGray] at (2.7,3.0) {LoRA adapters ($r{=}8$, Q/V only):};
  \node[lorachip, draw=cRGB] at (2.7,2.55) {LoRA$_{\mathrm{rgb}}$};
  \node[lorachip, draw=cDep] at (2.7,2.05) {LoRA$_{\mathrm{depth}}$};
  \node[lorachip, draw=cEvt] at (2.7,1.55) {LoRA$_{\mathrm{event}}$};
  \node[lorachip, draw=cLid] at (2.7,1.05) {LoRA$_{\mathrm{lidar}}$};
  \foreach \i in {1,2,3,4}{ \draw[flow] (in\i.east) -- (in\i.east -| enc.west); }
  \node[font=\tiny, text=cGray, align=center] at (2.7,0.2)
    {one pass per modality;\\identity = active LoRA};

  % ================= feature stack F_m =================
  \foreach \i/\c in {3/cLid,2/cEvt,1/cDep,0/cRGB}
    \node[card, draw=\c] (F\i) at (5.35+\i*0.09, 2.45+\i*0.09) {};
  \node[font=\scriptsize] at (5.5,1.8) {$F_m$};
  \node[font=\tiny, text=cGray] at (5.5,1.45) {$(1024,64^2)$};
  \draw[flow] (enc.east) -- (5.0,2.55);

  % ================= cross-modal attention (top row) =================
  \node[box, minimum width=26mm, anchor=west] (attn) at (6.4,3.3)
    {cross-modal attn $\times 2$\\ \tiny Q: own tokens; K,V: others};
  \draw[flow] (F3.north) |- ($(attn.west)+(0,0.05)$);
  % struck-out removed bias
  \node[font=\tiny, text=cGray, align=center, anchor=north] (nobias)
    at ($(attn.south)+(0,-0.08)$) {reliability logit bias};
  \draw[cGray, thick] ($(nobias.west)+(0.05,0)$) -- ($(nobias.east)+(-0.05,0)$);
  \node[font=\tiny, text=cGray, anchor=north, inner sep=1pt]
    at (nobias.south) {removed (negative result)};

  % ================= competence gate =================
  \node[box, minimum width=27mm, anchor=west] (gate) at (9.6,3.3)
    {competence gate\\ $\mathrm{softmax}_m(r_m/\tau)$ + veto floor};
  \draw[flow] (attn.east) -- (gate.west) node[midway, above, font=\tiny] {$\tilde F_m$};

  % ================= FPN + head =================
  \node[box, anchor=west] (fpn) at (13.7,3.3) {Simple\\FPN};
  \draw[flow] (gate.east) -- (fpn.west)
    node[midway, below, font=\tiny] {$F^{\mathrm{fused}}$};
  \node[box, anchor=west] (head) at (15.4,3.3) {seg\\head};
  \draw[flow] (fpn.east) -- (head.west);
  \node[opn] (plus) at (17.0,3.3) {$+$};
  \draw[flow] (head.east) -- (plus.west);
  \node[font=\normalsize, anchor=west] (yhat) at (17.5,3.3) {$\hat y$};
  \draw[flow] (plus.east) -- (yhat.west);

  % ================= bottom row: aux decoders + signals =================
  \node[box, minimum width=20mm, anchor=west] (aux) at (6.4,0.9)
    {aux decoders\\$D_m\times 4$};
  \draw[flow] (F0.south) |- (aux.west);
  \foreach \i/\c in {3/cLid,2/cEvt,1/cDep,0/cRGB}
    \node[card, draw=\c, minimum width=5mm, minimum height=4mm]
      (z\i) at (9.05+\i*0.08, 0.85+\i*0.08) {};
  \node[font=\scriptsize] at (9.15,0.35) {$z_m$};
  \draw[flow] (aux.east) -- (8.75,0.9);
  \node[box, draw=cRel, text=cRel, minimum width=24mm, anchor=west]
    (relsig) at (9.9,0.9)
    {reliability signals\\ $r_m$ (calibrated), $v_m$ (veto)};
  \draw[flow] (9.55,0.9) -- (relsig.west);
  \draw[sig] (relsig.north) |- ($(gate.south)+(0,-0.15)$) -- (gate.south)
    node[pos=0.15, right, font=\tiny, text=cRel] {$r_m, v_m$};

  % ================= router =================
  \node[box, draw=OursBlue, text=OursBlue, minimum width=30mm, anchor=west]
    (router) at (13.4,0.9)
    {\textbf{per-class reliability-}\\\textbf{anchored router (PRR)}};
  \draw[sig] (relsig.east) -- (router.west)
    node[midway, above, font=\tiny, text=cRel] {$\mathrm{sg}[r_m]$};
  \draw[flow] (z0.south) -- (9.05,0.15) -- (14.0,0.15) -- (14.0,0.45)
    node[pos=0.5, above, font=\tiny] {$z_m$};
  % thin F_m bus into the router
  \draw[->, cGray, thin] (F0.east) -- (6.1,2.45) -- (6.1,1.7) -- (12.9,1.7)
    -- (12.9,1.35) -- ($(router.north)+(-0.6,0)$)
    node[pos=0.5, above, font=\tiny, text=cGray] {$F_m$};
  \draw[flow, OursBlue] (router.east) -| (plus.south)
    node[pos=0.9, right, font=\tiny, text=OursBlue]
    {$\alpha\, z^{\mathrm{route}}$ ($\alpha$ zero-init)};

  % ================= training-only losses =================
  \node[loss, anchor=north] (laux) at (7.4,-0.15)
    {$\mathcal{L}_{\mathrm{aux}}$, $\mathcal{L}_{\mathrm{cal}}$};
  \draw[<-, cGray, dashed] (aux.south) -- (laux.north);
  \node[loss, anchor=north] (lroute) at (14.9,-0.15) {$\mathcal{L}_{\mathrm{route}}$};
  \draw[<-, cGray, dashed] (router.south) -- (lroute.north);
  \node[loss, anchor=south] (lohem) at (17.6,3.9) {$\mathcal{L}_{\mathrm{OHEM}}$};
  \draw[<-, cGray, dashed] (yhat.north) -- (lohem.south);

  % ================= legend =================
  \node[font=\tiny, anchor=west, align=left] at (0.0,-0.5)
    {\tikz{\draw[->, thick] (0,0)--(0.35,0);}\; inference graph\quad
     \tikz{\draw[->, cGray, dashed] (0,0)--(0.35,0);}\; {\color{cGray}training-only}\quad
     \tikz{\draw[->, cRel, thick] (0,0)--(0.35,0);}\; {\color{cRel}reliability signals (detached)}};
\end{tikzpicture}
}
% Structure source: notes/fact_method.md §0–§9 (no numbers to update).
% Suggested caption (place in the figure* env):
%   \caption{ReliaDINO. All four modalities are encoded by a single frozen
%   DINOv3 ViT-L/16; modality identity enters only through independent
%   low-rank adapters (color-coded, rank~8 on Q/V of all 24 blocks).
%   Per-modality auxiliary decoders provide the class posteriors from which
%   calibrated reliability $r_m$ and the veto signal $v_m$ are computed
%   (purple; detached). Modalities exchange information through two shared
%   cross-modal attention layers; the fused map is produced by a calibrated
%   competence gate with a training-free veto floor, decoded by a SimpleFPN
%   head. The per-class reliability-anchored router adds a residual routed
%   prediction scaled by a zero-initialized $\alpha$. Reliability acts only
%   at the gate and router: the attention-logit bias site is removed
%   (struck out; Sec.~\todo{ref}). Dashed gray parts are training-only.}
% ============================================================
\definecolor{cRGB}{RGB}{213,94,0}
\definecolor{cDep}{RGB}{0,114,178}
\definecolor{cEvt}{RGB}{230,159,0}
\definecolor{cLid}{RGB}{0,158,115}
\definecolor{cRel}{RGB}{136,78,160}
\definecolor{cGray}{RGB}{120,120,120}
\definecolor{OursBlue}{RGB}{0,90,181}
\begin{tikzpicture}[font=\scriptsize, >={Stealth[length=1.5mm]}, line cap=round,
  box/.style={draw, rounded corners=2pt, align=center, inner sep=3pt,
              minimum height=9mm},
  mod/.style={draw, very thick, minimum width=7.5mm, minimum height=6mm,
              inner sep=1pt, align=center, font=\tiny},
  lorachip/.style={draw, very thick, rounded corners=1.5pt, minimum width=14mm,
              minimum height=3.6mm, inner sep=1pt, font=\tiny, fill=white},
  card/.style={draw, fill=white, minimum width=6.5mm, minimum height=5mm,
              inner sep=0pt},
  loss/.style={draw, dashed, rounded corners=2pt, text=cGray, draw=cGray,
              align=center, inner sep=2.5pt, font=\tiny},
  sig/.style={->, cRel, thick},
  flow/.style={->, thick},
  opn/.style={circle, draw, thick, inner sep=0.5pt, font=\footnotesize}]

  % ================= inputs =================
  \node[mod, draw=cRGB] (in1) at (0,3.9)  {RGB};
  \node[mod, draw=cDep] (in2) at (0,3.0)  {Depth};
  \node[mod, draw=cEvt] (in3) at (0,2.1)  {Event};
  \node[mod, draw=cLid] (in4) at (0,1.2)  {LiDAR};
  \node[font=\tiny, text=cGray, align=center] at (0,0.55) {$x_m$: $4{\times}(3,1024^2)$};

  % ================= shared frozen encoder =================
  \node[box, minimum width=32mm, minimum height=34mm, anchor=west]
    (enc) at (1.1,2.55) {};
  \node[align=center, anchor=north] at (2.7,4.15)
    {\textbf{Frozen DINOv3 ViT-L/16}\\shared, 24 blocks};
  \node[font=\tiny, text=cGray] at (2.7,3.0) {LoRA adapters ($r{=}8$, Q/V only):};
  \node[lorachip, draw=cRGB] at (2.7,2.55) {LoRA$_{\mathrm{rgb}}$};
  \node[lorachip, draw=cDep] at (2.7,2.05) {LoRA$_{\mathrm{depth}}$};
  \node[lorachip, draw=cEvt] at (2.7,1.55) {LoRA$_{\mathrm{event}}$};
  \node[lorachip, draw=cLid] at (2.7,1.05) {LoRA$_{\mathrm{lidar}}$};
  \foreach \i in {1,2,3,4}{ \draw[flow] (in\i.east) -- (in\i.east -| enc.west); }
  \node[font=\tiny, text=cGray, align=center] at (2.7,0.2)
    {one pass per modality;\\identity = active LoRA};

  % ================= feature stack F_m =================
  \foreach \i/\c in {3/cLid,2/cEvt,1/cDep,0/cRGB}
    \node[card, draw=\c] (F\i) at (5.35+\i*0.09, 2.45+\i*0.09) {};
  \node[font=\scriptsize] at (5.5,1.8) {$F_m$};
  \node[font=\tiny, text=cGray] at (5.5,1.45) {$(1024,64^2)$};
  \draw[flow] (enc.east) -- (5.0,2.55);

  % ================= cross-modal attention (top row) =================
  \node[box, minimum width=26mm, anchor=west] (attn) at (6.4,3.3)
    {cross-modal attn $\times 2$\\ \tiny Q: own tokens; K,V: others};
  \draw[flow] (F3.north) |- ($(attn.west)+(0,0.05)$);
  % struck-out removed bias
  \node[font=\tiny, text=cGray, align=center, anchor=north] (nobias)
    at ($(attn.south)+(0,-0.08)$) {reliability logit bias};
  \draw[cGray, thick] ($(nobias.west)+(0.05,0)$) -- ($(nobias.east)+(-0.05,0)$);
  \node[font=\tiny, text=cGray, anchor=north, inner sep=1pt]
    at (nobias.south) {removed (negative result)};

  % ================= competence gate =================
  \node[box, minimum width=27mm, anchor=west] (gate) at (9.6,3.3)
    {competence gate\\ $\mathrm{softmax}_m(r_m/\tau)$ + veto floor};
  \draw[flow] (attn.east) -- (gate.west) node[midway, above, font=\tiny] {$\tilde F_m$};

  % ================= FPN + head =================
  \node[box, anchor=west] (fpn) at (13.7,3.3) {Simple\\FPN};
  \draw[flow] (gate.east) -- (fpn.west)
    node[midway, below, font=\tiny] {$F^{\mathrm{fused}}$};
  \node[box, anchor=west] (head) at (15.4,3.3) {seg\\head};
  \draw[flow] (fpn.east) -- (head.west);
  \node[opn] (plus) at (17.0,3.3) {$+$};
  \draw[flow] (head.east) -- (plus.west);
  \node[font=\normalsize, anchor=west] (yhat) at (17.5,3.3) {$\hat y$};
  \draw[flow] (plus.east) -- (yhat.west);

  % ================= bottom row: aux decoders + signals =================
  \node[box, minimum width=20mm, anchor=west] (aux) at (6.4,0.9)
    {aux decoders\\$D_m\times 4$};
  \draw[flow] (F0.south) |- (aux.west);
  \foreach \i/\c in {3/cLid,2/cEvt,1/cDep,0/cRGB}
    \node[card, draw=\c, minimum width=5mm, minimum height=4mm]
      (z\i) at (9.05+\i*0.08, 0.85+\i*0.08) {};
  \node[font=\scriptsize] at (9.15,0.35) {$z_m$};
  \draw[flow] (aux.east) -- (8.75,0.9);
  \node[box, draw=cRel, text=cRel, minimum width=24mm, anchor=west]
    (relsig) at (9.9,0.9)
    {reliability signals\\ $r_m$ (calibrated), $v_m$ (veto)};
  \draw[flow] (9.55,0.9) -- (relsig.west);
  \draw[sig] (relsig.north) |- ($(gate.south)+(0,-0.15)$) -- (gate.south)
    node[pos=0.15, right, font=\tiny, text=cRel] {$r_m, v_m$};

  % ================= router =================
  \node[box, draw=OursBlue, text=OursBlue, minimum width=30mm, anchor=west]
    (router) at (13.4,0.9)
    {\textbf{per-class reliability-}\\\textbf{anchored router (PRR)}};
  \draw[sig] (relsig.east) -- (router.west)
    node[midway, above, font=\tiny, text=cRel] {$\mathrm{sg}[r_m]$};
  \draw[flow] (z0.south) -- (9.05,0.15) -- (14.0,0.15) -- (14.0,0.45)
    node[pos=0.5, above, font=\tiny] {$z_m$};
  % thin F_m bus into the router
  \draw[->, cGray, thin] (F0.east) -- (6.1,2.45) -- (6.1,1.7) -- (12.9,1.7)
    -- (12.9,1.35) -- ($(router.north)+(-0.6,0)$)
    node[pos=0.5, above, font=\tiny, text=cGray] {$F_m$};
  \draw[flow, OursBlue] (router.east) -| (plus.south)
    node[pos=0.9, right, font=\tiny, text=OursBlue]
    {$\alpha\, z^{\mathrm{route}}$ ($\alpha$ zero-init)};

  % ================= training-only losses =================
  \node[loss, anchor=north] (laux) at (7.4,-0.15)
    {$\mathcal{L}_{\mathrm{aux}}$, $\mathcal{L}_{\mathrm{cal}}$};
  \draw[<-, cGray, dashed] (aux.south) -- (laux.north);
  \node[loss, anchor=north] (lroute) at (14.9,-0.15) {$\mathcal{L}_{\mathrm{route}}$};
  \draw[<-, cGray, dashed] (router.south) -- (lroute.north);
  \node[loss, anchor=south] (lohem) at (17.6,3.9) {$\mathcal{L}_{\mathrm{OHEM}}$};
  \draw[<-, cGray, dashed] (yhat.north) -- (lohem.south);

  % ================= legend =================
  \node[font=\tiny, anchor=west, align=left] at (0.0,-0.5)
    {\tikz{\draw[->, thick] (0,0)--(0.35,0);}\; inference graph\quad
     \tikz{\draw[->, cGray, dashed] (0,0)--(0.35,0);}\; {\color{cGray}training-only}\quad
     \tikz{\draw[->, cRel, thick] (0,0)--(0.35,0);}\; {\color{cRel}reliability signals (detached)}};
\end{tikzpicture}
}
% Structure source: notes/fact_method.md §0–§9 (no numbers to update).
% Suggested caption (place in the figure* env):
%   \caption{ReliaDINO. All four modalities are encoded by a single frozen
%   DINOv3 ViT-L/16; modality identity enters only through independent
%   low-rank adapters (color-coded, rank~8 on Q/V of all 24 blocks).
%   Per-modality auxiliary decoders provide the class posteriors from which
%   calibrated reliability $r_m$ and the veto signal $v_m$ are computed
%   (purple; detached). Modalities exchange information through two shared
%   cross-modal attention layers; the fused map is produced by a calibrated
%   competence gate with a training-free veto floor, decoded by a SimpleFPN
%   head. The per-class reliability-anchored router adds a residual routed
%   prediction scaled by a zero-initialized $\alpha$. Reliability acts only
%   at the gate and router: the attention-logit bias site is removed
%   (struck out; Sec.~\todo{ref}). Dashed gray parts are training-only.}
% ============================================================
\definecolor{cRGB}{RGB}{213,94,0}
\definecolor{cDep}{RGB}{0,114,178}
\definecolor{cEvt}{RGB}{230,159,0}
\definecolor{cLid}{RGB}{0,158,115}
\definecolor{cRel}{RGB}{136,78,160}
\definecolor{cGray}{RGB}{120,120,120}
\definecolor{OursBlue}{RGB}{0,90,181}
\begin{tikzpicture}[font=\scriptsize, >={Stealth[length=1.5mm]}, line cap=round,
  box/.style={draw, rounded corners=2pt, align=center, inner sep=3pt,
              minimum height=9mm},
  mod/.style={draw, very thick, minimum width=7.5mm, minimum height=6mm,
              inner sep=1pt, align=center, font=\tiny},
  lorachip/.style={draw, very thick, rounded corners=1.5pt, minimum width=14mm,
              minimum height=3.6mm, inner sep=1pt, font=\tiny, fill=white},
  card/.style={draw, fill=white, minimum width=6.5mm, minimum height=5mm,
              inner sep=0pt},
  loss/.style={draw, dashed, rounded corners=2pt, text=cGray, draw=cGray,
              align=center, inner sep=2.5pt, font=\tiny},
  sig/.style={->, cRel, thick},
  flow/.style={->, thick},
  opn/.style={circle, draw, thick, inner sep=0.5pt, font=\footnotesize}]

  % ================= inputs =================
  \node[mod, draw=cRGB] (in1) at (0,3.9)  {RGB};
  \node[mod, draw=cDep] (in2) at (0,3.0)  {Depth};
  \node[mod, draw=cEvt] (in3) at (0,2.1)  {Event};
  \node[mod, draw=cLid] (in4) at (0,1.2)  {LiDAR};
  \node[font=\tiny, text=cGray, align=center] at (0,0.55) {$x_m$: $4{\times}(3,1024^2)$};

  % ================= shared frozen encoder =================
  \node[box, minimum width=32mm, minimum height=34mm, anchor=west]
    (enc) at (1.1,2.55) {};
  \node[align=center, anchor=north] at (2.7,4.15)
    {\textbf{Frozen DINOv3 ViT-L/16}\\shared, 24 blocks};
  \node[font=\tiny, text=cGray] at (2.7,3.0) {LoRA adapters ($r{=}8$, Q/V only):};
  \node[lorachip, draw=cRGB] at (2.7,2.55) {LoRA$_{\mathrm{rgb}}$};
  \node[lorachip, draw=cDep] at (2.7,2.05) {LoRA$_{\mathrm{depth}}$};
  \node[lorachip, draw=cEvt] at (2.7,1.55) {LoRA$_{\mathrm{event}}$};
  \node[lorachip, draw=cLid] at (2.7,1.05) {LoRA$_{\mathrm{lidar}}$};
  \foreach \i in {1,2,3,4}{ \draw[flow] (in\i.east) -- (in\i.east -| enc.west); }
  \node[font=\tiny, text=cGray, align=center] at (2.7,0.2)
    {one pass per modality;\\identity = active LoRA};

  % ================= feature stack F_m =================
  \foreach \i/\c in {3/cLid,2/cEvt,1/cDep,0/cRGB}
    \node[card, draw=\c] (F\i) at (5.35+\i*0.09, 2.45+\i*0.09) {};
  \node[font=\scriptsize] at (5.5,1.8) {$F_m$};
  \node[font=\tiny, text=cGray] at (5.5,1.45) {$(1024,64^2)$};
  \draw[flow] (enc.east) -- (5.0,2.55);

  % ================= cross-modal attention (top row) =================
  \node[box, minimum width=26mm, anchor=west] (attn) at (6.4,3.3)
    {cross-modal attn $\times 2$\\ \tiny Q: own tokens; K,V: others};
  \draw[flow] (F3.north) |- ($(attn.west)+(0,0.05)$);
  % struck-out removed bias
  \node[font=\tiny, text=cGray, align=center, anchor=north] (nobias)
    at ($(attn.south)+(0,-0.08)$) {reliability logit bias};
  \draw[cGray, thick] ($(nobias.west)+(0.05,0)$) -- ($(nobias.east)+(-0.05,0)$);
  \node[font=\tiny, text=cGray, anchor=north, inner sep=1pt]
    at (nobias.south) {removed (negative result)};

  % ================= competence gate =================
  \node[box, minimum width=27mm, anchor=west] (gate) at (9.6,3.3)
    {competence gate\\ $\mathrm{softmax}_m(r_m/\tau)$ + veto floor};
  \draw[flow] (attn.east) -- (gate.west) node[midway, above, font=\tiny] {$\tilde F_m$};

  % ================= FPN + head =================
  \node[box, anchor=west] (fpn) at (13.7,3.3) {Simple\\FPN};
  \draw[flow] (gate.east) -- (fpn.west)
    node[midway, below, font=\tiny] {$F^{\mathrm{fused}}$};
  \node[box, anchor=west] (head) at (15.4,3.3) {seg\\head};
  \draw[flow] (fpn.east) -- (head.west);
  \node[opn] (plus) at (17.0,3.3) {$+$};
  \draw[flow] (head.east) -- (plus.west);
  \node[font=\normalsize, anchor=west] (yhat) at (17.5,3.3) {$\hat y$};
  \draw[flow] (plus.east) -- (yhat.west);

  % ================= bottom row: aux decoders + signals =================
  \node[box, minimum width=20mm, anchor=west] (aux) at (6.4,0.9)
    {aux decoders\\$D_m\times 4$};
  \draw[flow] (F0.south) |- (aux.west);
  \foreach \i/\c in {3/cLid,2/cEvt,1/cDep,0/cRGB}
    \node[card, draw=\c, minimum width=5mm, minimum height=4mm]
      (z\i) at (9.05+\i*0.08, 0.85+\i*0.08) {};
  \node[font=\scriptsize] at (9.15,0.35) {$z_m$};
  \draw[flow] (aux.east) -- (8.75,0.9);
  \node[box, draw=cRel, text=cRel, minimum width=24mm, anchor=west]
    (relsig) at (9.9,0.9)
    {reliability signals\\ $r_m$ (calibrated), $v_m$ (veto)};
  \draw[flow] (9.55,0.9) -- (relsig.west);
  \draw[sig] (relsig.north) |- ($(gate.south)+(0,-0.15)$) -- (gate.south)
    node[pos=0.15, right, font=\tiny, text=cRel] {$r_m, v_m$};

  % ================= router =================
  \node[box, draw=OursBlue, text=OursBlue, minimum width=30mm, anchor=west]
    (router) at (13.4,0.9)
    {\textbf{per-class reliability-}\\\textbf{anchored router (PRR)}};
  \draw[sig] (relsig.east) -- (router.west)
    node[midway, above, font=\tiny, text=cRel] {$\mathrm{sg}[r_m]$};
  \draw[flow] (z0.south) -- (9.05,0.15) -- (14.0,0.15) -- (14.0,0.45)
    node[pos=0.5, above, font=\tiny] {$z_m$};
  % thin F_m bus into the router
  \draw[->, cGray, thin] (F0.east) -- (6.1,2.45) -- (6.1,1.7) -- (12.9,1.7)
    -- (12.9,1.35) -- ($(router.north)+(-0.6,0)$)
    node[pos=0.5, above, font=\tiny, text=cGray] {$F_m$};
  \draw[flow, OursBlue] (router.east) -| (plus.south)
    node[pos=0.9, right, font=\tiny, text=OursBlue]
    {$\alpha\, z^{\mathrm{route}}$ ($\alpha$ zero-init)};

  % ================= training-only losses =================
  \node[loss, anchor=north] (laux) at (7.4,-0.15)
    {$\mathcal{L}_{\mathrm{aux}}$, $\mathcal{L}_{\mathrm{cal}}$};
  \draw[<-, cGray, dashed] (aux.south) -- (laux.north);
  \node[loss, anchor=north] (lroute) at (14.9,-0.15) {$\mathcal{L}_{\mathrm{route}}$};
  \draw[<-, cGray, dashed] (router.south) -- (lroute.north);
  \node[loss, anchor=south] (lohem) at (17.6,3.9) {$\mathcal{L}_{\mathrm{OHEM}}$};
  \draw[<-, cGray, dashed] (yhat.north) -- (lohem.south);

  % ================= legend =================
  \node[font=\tiny, anchor=west, align=left] at (0.0,-0.5)
    {\tikz{\draw[->, thick] (0,0)--(0.35,0);}\; inference graph\quad
     \tikz{\draw[->, cGray, dashed] (0,0)--(0.35,0);}\; {\color{cGray}training-only}\quad
     \tikz{\draw[->, cRel, thick] (0,0)--(0.35,0);}\; {\color{cRel}reliability signals (detached)}};
\end{tikzpicture}
}
  \caption{ReliaDINO. All four modalities are encoded by a single frozen
  DINOv3 ViT-L/16; modality identity enters only through independent
  low-rank adapters (color-coded, rank~8 on Q/V of all 24 blocks).
  Per-modality auxiliary decoders provide the class posteriors from which
  calibrated reliability $r_m$ and the veto signal $v_m$ are computed
  (purple; detached). Modalities exchange information through two shared
  cross-modal attention layers; the fused map is produced by a calibrated
  competence gate with a training-free veto floor, decoded by a SimpleFPN
  head. The per-class reliability-anchored router adds a residual routed
  prediction scaled by a zero-initialized $\alpha$. Reliability acts only
  at the gate and router: the attention-logit bias site is removed
  (struck out; Sec.~\ref{sec:method:negative}). Dashed gray parts are
  training-only.}
  \label{fig:arch}
\end{figure*}

% ============================================================
% Section 3 — Method (body-only; \input from root.tex)
% Model naming: ReliaDINO (full = with PRR router);
% "ReliaDINO (no router)" = ablation variant. No internal codes.
% All numbers from notes/fact_method.md; missing items marked \todo{}.
% ============================================================

\section{Method}
\label{sec:method}

\subsection{Overview}
\label{sec:method:overview}

Fig.~\ref{fig:arch} gives an overview of ReliaDINO. Let
$\mathcal{M}=\{\mathrm{rgb},\mathrm{depth},\mathrm{event},\mathrm{lidar}\}$
denote the modality set with $M=|\mathcal{M}|=4$, and let $C$ be the number of
semantic classes ($C{=}25$ on DELIVER). Each modality is rendered as a
three-channel image $x_m\in\mathbb{R}^{B\times 3\times H\times W}$
($H{=}W{=}1024$) and encoded by a \emph{single shared frozen} DINOv3
ViT-L/16 backbone~\cite{dinov3_2025} $E(\cdot;\theta_0)$, whose attention
projections carry independent per-modality low-rank adapters
$\Delta\theta_m$ (Sec.~\ref{sec:method:lora}). This yields one stride-16
token map $F_m\in\mathbb{R}^{B\times d\times h\times w}$ per modality
($d{=}1024$, $h{=}w{=}64$). A light auxiliary decoder $D_m$ produces
per-modality class logits $z_m=D_m(F_m)$, from which all reliability signals
are derived: a calibrated self-entropy reliability $r_m$, and a
corroboration-based veto signal $v_m$ (Sec.~\ref{sec:method:fusion}).
The modalities exchange information through two shared cross-modal attention
layers, producing attended maps $\tilde F_m$, which are fused by a calibrated
competence gate with weights $w^{g}_m$ and a training-free veto floor. The
fused map passes through a ViTDet-style feature
pyramid\todo{cite ViTDet} and a simple FPN segmentation head. Finally, the
Per-class Reliability-anchored Router (PRR, Sec.~\ref{sec:method:router};
Fig.~\ref{fig:router}) computes per-class, per-pixel modality weights
$w^{r}$ that mix the per-modality logits $z_m$ into a routed prediction,
added to the head output as a zero-initialized residual:
\begin{equation}
\hat y \;=\; \mathrm{Head}\big(\mathrm{FPN}(F^{\mathrm{fused}})\big)
\;+\; \alpha\cdot \mathrm{up}_{\times 4}\!\big(z^{\mathrm{route}}\big),
\label{eq:overview}
\end{equation}
where $\alpha$ is a learnable scalar initialized to zero. Reliability thus
enters the model at exactly two places---the output-level competence gate and
the per-class router---and, deliberately, \emph{not} inside the attention
logits (Sec.~\ref{sec:method:negative}).

\subsection{Frozen Backbone with Per-Modality Low-Rank Adaptation}
\label{sec:method:lora}

All $\theta_0$ backbone parameters remain frozen; modality identity enters the
encoder \emph{only} through the active LoRA~\cite{lora2022} adapter, which is
switched per forward pass ($M$ sequential forwards per step). For every one of
the 24 attention blocks, LoRA is injected on the fused query--key--value
projection $W^{qkv}_{\ell}$, restricted to the query and value output slices
(the key slice is untouched), following the LoRA-on-promptable-segmenter
convention~\cite{samed2023}. For block $\ell$ and active modality $m$, with
token input $x\in\mathbb{R}^{d}$,
\begin{equation}
\begin{aligned}
{}[q;\,k;\,v] &= W^{qkv}_{\ell}x + b^{qkv}_{\ell},\\
q &\mathrel{+}= \tfrac{\alpha_L}{r}\, B^{q}_{m,\ell} A^{q}_{m,\ell}\, x,
\qquad
v \mathrel{+}= \tfrac{\alpha_L}{r}\, B^{v}_{m,\ell} A^{v}_{m,\ell}\, x,
\end{aligned}
\label{eq:lora}
\end{equation}
with $A\in\mathbb{R}^{r\times d}$, $B\in\mathbb{R}^{d\times r}$, rank $r=8$,
and $\alpha_L=r$ so that the effective scale is $1$. The $B$ matrices are
zero-initialized, so every modality branch starts as the identity mapping of
the pretrained backbone. This design \emph{multiplexes} one visual foundation
model across four heterogeneous sensors: each modality receives its own
low-rank subspace for input-statistics adaptation, while the frozen weights
prevent catastrophic forgetting of the pretrained representation. The
adapters total $4\times r\,d$ parameters per (block, modality), i.e.\
$\approx$3.1\,M parameters for all $24\times 4$ adapters---roughly $1\%$ of
the frozen backbone's parameters. The no-router variant totals 349.9\,M
parameters of which 46.8\,M (13.4\%) are trainable (adapters, fusion,
decoders, head); the router adds only $\approx$0.27\,M
(\todo{exact full-model parameter count from its training log}).

\subsection{Cross-Modal Attention Fusion}
\label{sec:method:fusion}

\subsubsection{Cross-modal attention}
Tokens of modality $m$, $t_m\in\mathbb{R}^{B\times N\times d}$ ($N{=}hw$),
attend to the concatenated tokens of the other $M{-}1$ modalities through two
pre-norm attention blocks whose weights are shared across modalities
($n_h{=}8$ heads, MLP ratio $4$):
\begin{equation}
\mathrm{Attn}(t_m) = \mathrm{softmax}\!\Big(\tfrac{QK^{\top}}{\sqrt{d/n_h}}\Big)V,
\label{eq:xattn}
\end{equation}
\begin{equation}
t_m \leftarrow t_m + W_{\mathrm{proj}}\,\mathrm{Attn}(t_m),\quad
t_m \leftarrow t_m + \mathrm{MLP}\big(\mathrm{LN}(t_m)\big),
\end{equation}
where $Q$ derives from $t_m$ and $K,V$ from
$\mathrm{concat}\{t_j\}_{j\neq m}$. This generalizes memory
attention over video frames~\cite{sam2_2024,memorysam2025} to
``frames''\,$=$\,modalities, but---unlike our earlier
designs---the attention in Eq.~\eqref{eq:xattn} is \emph{pure}: no
reliability bias is added to its logits (Sec.~\ref{sec:method:negative}),
which also permits fused attention kernels. The output tokens are reshaped
back to per-modality maps $\tilde F_m$.

\subsubsection{Calibrated reliability}
Each auxiliary decoder ($\mathrm{Conv}3{\times}3\!\to\!\mathrm{GN}\!\to\!
\mathrm{GELU}\!\to\!\mathrm{Conv}1{\times}1$) predicts per-modality class
logits $z_m$ at token resolution. With a learned per-modality temperature
$T_m=\exp(\theta^{T}_m)$ (clamped to $[0.05,20]$, initialized to $1$), the
calibrated self-entropy reliability is
\begin{equation}
\hat p_m = \mathrm{softmax}(z_m/T_m),\qquad
r_m = 1-\frac{H(\hat p_m)}{\log C},
\label{eq:relcal}
\end{equation}
with $H(\hat p)=-\sum_{k=1}^{C}\hat p_k\log\hat p_k$; $r_m$ is a per-pixel
map in $[0,1]$. All signals are computed from detached logits, so the
decision path does not back-propagate into the decoders through them.

\subsubsection{Competence gate with veto floor}
The attended maps are fused by a softmax gate over modalities,
\begin{equation}
w^{g} = \mathrm{softmax}_{m}\!\big(r_m/\tau\big),\qquad \tau=0.25,
\label{eq:gate}
\end{equation}
computed pixel-wise. The gate consumes only the \emph{calibrated}
self-reliability $r_m$; cross-modal corroboration is intentionally excluded
here, since agreement with the consensus mis-ranks degraded sensors as
reliable in easy regions. Corroboration instead drives a training-free
\emph{veto}: with uncalibrated confidences $s_m=1-H(\mathrm{softmax}(z_m))/\log C$,
leave-one-out consensus $c_m=\frac{1}{M-1}\big(\sum_j p_j-p_m\big)_{+}$, and
Bhattacharyya corroboration
$\mathrm{corr}_m=\sum_{k}\sqrt{p_{m,k}\,c_{m,k}}$, the veto signal blends in
a unique-confidence guard $g_m=\mathrm{clamp}(s_m-\max_{j\neq m}s_j,0,1)$
that protects a modality that is alone-confident where the consensus is
blind:
\begin{equation}
v_m = \mathrm{clamp}\big(g_m s_m + (1-g_m)\,\mathrm{corr}_m,\ 0,\ 1\big).
\label{eq:veto}
\end{equation}
Wherever $v_m<\delta$ (the modality is both consensus-contradicted and
self-unconfident) its gate weight is capped and the weights renormalized:
\begin{equation}
w^{g}_m \leftarrow \min\!\big(w^{g}_m,\kappa\big)\ \text{if}\ v_m<\delta,
\qquad
w^{g} \leftarrow w^{g}\big/\textstyle\sum_m w^{g}_m,
\label{eq:vetofloor}
\end{equation}
with $\delta=0.10$, $\kappa=0.05$. The fused representation is the gated sum
$F^{\mathrm{fused}}=\sum_m w^{g}_m\odot\tilde F_m$; the final configuration
uses this gated fusion (rather than an equal average), which contributes
$+0.26$ test mIoU at full resolution (Sec.~\ref{sec:experiments}). A single
feature pyramid then expands $F^{\mathrm{fused}}$ to strides
$\{4,8,16,32\}$, and a query-free head sums all levels at stride~4, applies
two $3{\times}3$ convolution blocks and a $1{\times}1$ classifier.

\subsection{Per-Class Reliability-Anchored Router}
\label{sec:method:router}

The scalar gate of Eq.~\eqref{eq:gate} assigns one weight per pixel to each
modality---the same weight for every class. This provably anti-selects for
classes whose best sensor is globally weak: e.g., at night we measure
per-modality competence for road lines of $0.798$ (RGB) versus $0.001$
(depth), yet the scalar gate still allocates $0.432$ of the mixture to depth.
The PRR (Fig.~\ref{fig:router}) therefore routes \emph{per class}: for each
class $k$ and pixel it produces modality weights
\begin{equation}
w^{r} = \mathrm{softmax}_{m}\!\Big(g_m(F_m) + \lambda_{\mathrm{anc}}\,
\mathrm{sg}[r_m]\Big)\in\mathbb{R}^{M\times B\times C\times h\times w},
\label{eq:router}
\end{equation}
with $\lambda_{\mathrm{anc}}=1$, where $g_m$ is a two-layer $1{\times}1$
convolutional head ($1024\!\to\!64\!\to\!C$) on the \emph{pre-fusion}
features, $\mathrm{sg}[\cdot]$ is the stop-gradient operator, and $r_m$
broadcasts over the $C$ class channels. The routed prediction is a mixture of
the per-modality logits, $z^{\mathrm{route}}=\sum_m w^{r}_m\odot z_m$, added
to the head output as in Eq.~\eqref{eq:overview}; the router is part of the
inference graph.

Two design choices make this router trainable where purely learned gates were
not. First, the last convolution of $g_m$ is \emph{zero-initialized} (weights
and bias) and $\alpha$ in Eq.~\eqref{eq:overview} starts at zero: at
initialization the model is exactly the no-router model and the routing, once
$\alpha$ grows, is initially driven purely by the calibrated reliability
anchor. In our experience, learned fusion gates trained from random
initialization on this task repeatedly collapsed to a constant,
input-independent mixture; anchoring the softmax logits to a training-free
reliability signal removes the flat-gradient start that causes this collapse,
while the learned term retains the capacity to correct the anchor per class.
Second, a \emph{decisive} regularizer replaces the naive diversity reward
(which provably pushes the router \emph{toward} uniformity): with per-pixel
mixing entropy
$H_{\mathrm{pix}}=\mathbb{E}_{b,k,s}\big[-\sum_m w^{r}_m\log w^{r}_m\big]$
and batch-marginal entropy $H_{\mathrm{bar}}$ computed on
$\bar w_m=\mathbb{E}_{b,s}[w^{r}_m]$ per class and then averaged,
\begin{equation}
\mathcal{L}_{\mathrm{route}} =
\lambda_{\mathrm{reg}}\big(H_{\mathrm{pix}} - H_{\mathrm{bar}}\big),
\qquad \lambda_{\mathrm{reg}}=0.01,
\label{eq:routereg}
\end{equation}
i.e., the router is rewarded for committing per pixel and class (low local
entropy) while keeping all modalities in use on average (high marginal
entropy), preventing global single-modality collapse.

\subsection{Training-Only Reliability Supervision}
\label{sec:method:training}

The auxiliary decoders are trained with two objectives that shape the
reliability signals; at inference they remain in the graph only as the signal
source for the gate and router (and as router mixture inputs). First, deep
supervision: each $z_m$ is bilinearly resized to quarter label resolution and
supervised with cross-entropy,
$\mathcal{L}_{\mathrm{auxCE}}=\frac{1}{M}\sum_m
\mathrm{CE}\big(\mathrm{up}(z_m),\,y{\downarrow_4}\big)$.
Second, a correctness-contrastive calibration loss drives the temperatures
$T_m$ and the decoders so that confidence \emph{ranks} correctness: with
normalized entropy $e_m = H(\mathrm{softmax}(z_m/T_m))/\log C$ and calibrated
prediction $\hat k=\arg\max_k \hat p_{m,k}$,
\begin{equation}
\mathcal{L}_{\mathrm{cal}} = \frac{1}{M}\sum_{m}\Big[
\mathbb{E}_{\hat k\neq y}\big[\,1-e_m\,\big]
+\mathbb{E}_{\hat k = y}\big[\,e_m\,\big]\Big],
\label{eq:calib}
\end{equation}
raising entropy on wrong pixels and lowering it on correct ones, which
directly optimizes the reliability signal's correctness-ranking (AUROC). The
total objective is
\begin{equation}
\mathcal{L} = \mathcal{L}_{\mathrm{OHEM}}(\hat y, y)
+ \lambda_{\mathrm{aux}}\,\mathcal{L}_{\mathrm{auxCE}}
+ \lambda_{\mathrm{cal}}\,\mathcal{L}_{\mathrm{cal}}
+ \mathcal{L}_{\mathrm{route}},
\label{eq:total}
\end{equation}
with online hard-example-mined cross-entropy~\cite{ohem2016} as the main
term, $\lambda_{\mathrm{aux}}=0.5$, and $\lambda_{\mathrm{cal}}=0.1$
($\mathcal{L}_{\mathrm{route}}$ only in the full model). No supervision
beyond the segmentation labels is used---no ground-truth depth, no condition
annotations, and no text encoders (contribution C3).

\subsection{Where Reliability Injection Fails}
\label{sec:method:negative}

An apparently natural alternative---and the design we pursued across five
model generations---injects reliability \emph{inside} the cross-modal
attention, as an additive pre-softmax bias on the key logits:
\begin{equation}
\mathrm{Attn}(t_m) = \mathrm{softmax}\!\Big(\tfrac{QK^{\top}}{\sqrt{d/n_h}}
+ \lambda_1 B^{\mathrm{cal}} + \lambda_2 B^{\mathrm{cons}}\Big)V,
\label{eq:attnbias}
\end{equation}
where a key token at location $s$ of modality $j$ carries the per-token
scalar $\lambda_1 B^{\mathrm{cal}}_j(s)+\lambda_2 B^{\mathrm{cons}}_j(s)$,
built from the modality-centered signals
$B^{\mathrm{cal}}_m=r_m-\frac{1}{M}\sum_j r_j$ and
$B^{\mathrm{cons}}_m=\mathrm{corr}_m-\frac{1}{M}\sum_j \mathrm{corr}_j$,
with learnable scales $\lambda_1,\lambda_2$. The intuition---down-weight
unreliable keys while preserving their values---is appealing, and
neighboring placements exist in the
literature~\cite{primed2026,sae2026,sam2long2024}. Yet on strong per-modality
LoRA features this bias is a no-op or a net harm: toggling it on trained
checkpoints changes test mIoU by $|\Delta|\le 0.03$ with fused-feature
cosine similarity of $1.0$, and on an earlier backbone generation the bias
was a statistically significant net loss ($p=4.5\times10^{-22}$). We
therefore \emph{remove} it from the final model, keeping reliability only at
the output gate and the per-class router, where class decisions are actually
formed. Sec.~\ref{sec:experiments} quantifies this negative result and its
interpretation: once per-modality adaptation aligns the modalities in feature
space, logit-level suppression inside attention has nothing left to fix.


% ---------------- Fig. 3: router detail (single-column) --------------------
\begin{figure}[t]
  \centering
  \resizebox{0.98\linewidth}{!}{% ============================================================
% fig:router — Per-class Reliability-anchored Router (PRR) detail
% BODY-ONLY: \input inside a single-column {figure}.
% Preamble needs: \usepackage{tikz, xcolor, amsmath, amssymb}
%   \usetikzlibrary{positioning, arrows.meta, calc}
% Structure source: notes/fact_method.md §6.
% Suggested caption (place in the figure env):
%   \caption{Per-class reliability-anchored router. A zero-initialized
%   per-modality head $g_m$ on the pre-fusion features is anchored by the
%   detached calibrated reliability $r_m$; a softmax over modalities yields
%   per-class, per-pixel routing weights that mix the per-modality
%   posteriors $z_m$ into a residual correction scaled by a
%   zero-initialized $\alpha$. At initialization the model is exactly the
%   no-router variant, and routing follows reliability alone (collapse-safe).
%   Motivation: the scalar competence gate anti-selects per class --- e.g.\
%   for road lines at night the per-modality competence is 0.80 (RGB) vs.\
%   0.00 (depth), yet the scalar gate assigns depth a weight of 0.43; a
%   per-class router lets each class route to the modality that sees it.}
% ============================================================
\definecolor{cRGB}{RGB}{213,94,0}
\definecolor{cDep}{RGB}{0,114,178}
\definecolor{cEvt}{RGB}{230,159,0}
\definecolor{cLid}{RGB}{0,158,115}
\definecolor{cRel}{RGB}{136,78,160}
\definecolor{cGray}{RGB}{120,120,120}
\definecolor{OursBlue}{RGB}{0,90,181}
\begin{tikzpicture}[font=\scriptsize, >={Stealth[length=1.5mm]}, line cap=round,
  box/.style={draw, rounded corners=2pt, align=center, inner sep=3pt,
              minimum height=7mm},
  card/.style={draw, fill=white, minimum width=5.5mm, minimum height=4.5mm,
              inner sep=0pt},
  note/.style={draw=cGray, fill=cGray!8, rounded corners=2pt, text=cGray,
              align=left, inner sep=3pt, font=\tiny},
  opn/.style={circle, draw, thick, inner sep=0.5pt, font=\footnotesize},
  flow/.style={->, thick},
  sig/.style={->, cRel, thick}]

  % ---------- inputs (labels above the tokens) ----------
  \foreach \i/\c in {3/cLid,2/cEvt,1/cDep,0/cRGB}
    \node[card, draw=\c] (F\i) at (0.6+\i*0.09, 6.0+\i*0.09) {};
  \node[align=center, font=\tiny] at (0.75,6.75) {$F_m$ (pre-fusion)\\$(1024,64^2)$};

  \node[card, draw=cRel, fill=cRel!12] (rm) at (2.8,6.05) {};
  \node[align=center, font=\tiny, text=cRel] at (2.8,6.75)
    {$r_m$ calibrated\\reliability $(1,64^2)$};

  \foreach \i/\c in {3/cLid,2/cEvt,1/cDep,0/cRGB}
    \node[card, draw=\c] (z\i) at (4.5+\i*0.09, 6.0+\i*0.09) {};
  \node[align=center, font=\tiny] at (4.65,6.75) {$z_m$ aux logits\\$(C,64^2)$};

  % ---------- router head ----------
  \node[box, minimum width=34mm] (gm) at (0.9,4.5)
    {$g_m$: Conv$1{\times}1$ $1024{\to}64$\\
     $\to$ ReLU $\to$ Conv$1{\times}1$ $64{\to}C$\\
     \tiny last conv \textbf{zero-init}};
  \draw[flow] (F0.south) -- (gm.north);

  % ---------- anchor add ----------
  \node[opn] (add1) at (2.4,3.4) {$+$};
  \draw[flow] (gm.south) |- (add1.west);
  \draw[sig] (rm.south) -- (2.8,3.4) -- (add1.east)
    node[pos=0.35, right, font=\tiny, text=cRel]
    {$\lambda_{\mathrm{anc}}\,\mathrm{sg}[r_m]$};

  % ---------- softmax over modalities ----------
  \node[box, minimum width=28mm] (smax) at (2.4,2.55)
    {softmax over modalities $m$};
  \draw[flow] (add1.south) -- (smax.north);
  \node[font=\tiny, anchor=west, text=OursBlue] at (smax.east)
    {\;$w^{r}\in\mathbb{R}^{M\times C\times h\times w}$};
  \node[font=\tiny, anchor=north, text=OursBlue] at (smax.south)
    {per-class, per-pixel modality weights};

  % ---------- routed mixture ----------
  \node[opn] (mul) at (2.4,1.35) {$\times$};
  \draw[flow] ($(smax.south)+(0,-0.3)$) -- (mul.north);
  \draw[flow] (z0.south) -- (4.5,1.35) -- (mul.east)
    node[pos=0.8, above, font=\tiny] {$z_m$};
  \node[font=\tiny, anchor=north west] at ($(mul.south)+(0.05,0)$) {$\sum_m$};
  \node[box] (zroute) at (0.7,1.35) {$z^{\mathrm{route}}$\\$(C,64^2)$};
  \draw[flow] (mul.west) -- (zroute.east);

  % ---------- residual add ----------
  \node[box, minimum width=16mm] (alpha) at (0.7,0.3)
    {$\times\,\alpha$\\ \tiny learnable, \textbf{zero-init}};
  \draw[flow] (zroute.south) -- (alpha.north);
  \node[opn] (add2) at (3.1,0.3) {$+$};
  \draw[flow] (alpha.east) -- (add2.west)
    node[midway, above, font=\tiny] {up$\times 4$};
  \node[box, draw=cGray, text=cGray] (headstub) at (3.1,-0.75)
    {Head(FPN($F^{\mathrm{fused}}$))};
  \draw[flow, cGray] (headstub.north) -- (add2.south);
  \node[font=\normalsize, anchor=west] (yhat) at (4.2,0.3) {$\hat y$};
  \draw[flow] (add2.east) -- (yhat.west);

  % ---------- side notes ----------
  \node[note, text width=25mm, anchor=north west] at (5.0,4.9)
    {\textbf{at init:} $g_m{\equiv}0$, $\alpha{=}0$\\
     $\Rightarrow$ $w^r=\mathrm{softmax}(\lambda_{\mathrm{anc}} r_m)$,\\
     output = no-router model\\ (collapse-safe start)};
  \node[note, text width=25mm, anchor=north west] at (5.0,2.9)
    {\textbf{decisive regularizer}\\
     $\mathcal{L}_{\mathrm{route}}=\lambda_{\mathrm{reg}}(H_{\mathrm{pix}}-H_{\mathrm{bar}})$\\
     commit per pixel/class,\\ stay diverse on average\\
     ($\lambda_{\mathrm{reg}}{=}0.01$)};
\end{tikzpicture}
}
  \caption{Per-class reliability-anchored router. A zero-initialized
  per-modality head $g_m$ on the pre-fusion features is anchored by the
  detached calibrated reliability $r_m$; a softmax over modalities yields
  per-class, per-pixel routing weights that mix the per-modality
  posteriors $z_m$ into a residual correction scaled by a
  zero-initialized $\alpha$. At initialization the model is exactly the
  no-router variant, and routing follows reliability alone
  (collapse-safe). Motivation: the scalar competence gate anti-selects per
  class---e.g.\ for road lines at night the per-modality competence is
  0.80 (RGB) vs.\ 0.00 (depth), yet the scalar gate assigns depth a weight
  of 0.43; a per-class router lets each class route to the modality that
  sees it.}
  \label{fig:router}
\end{figure}

% =====================================================================
% Section IV — Experiments (body-only; \input into root.tex)
% All numbers sourced from notes/fact_experiments.md (2026-07-15).
% Internal version codes must never appear here.
% =====================================================================
\section{Experiments}
\label{sec:experiments}

\subsection{Setup}
\label{sec:setup}

\textbf{Dataset and metric.}
We evaluate on DELIVER~\cite{cmnext2023}, a synthetic driving benchmark
with four spatially aligned modalities (RGB, depth, event, LiDAR),
25 semantic classes, and 3{,}983 / 2{,}005 / 1{,}897 images for
training / validation / test at $1042\times1042$ resolution. The data
cover five environmental conditions (cloudy, foggy, night, rainy,
sunny) as well as partial sensor-failure corruption cases. We report
mean Intersection-over-Union (mIoU) on both the validation and test
splits, and use the four-modality (camera--LiDAR--depth--event)
protocol throughout. MUSES~\cite{muses2024} results are planned:
\todo{MUSES official-protocol numbers (test-server submission pending;
only a non-comparable internal letterbox evaluation exists)}.

\textbf{Implementation details.}
The backbone is a single frozen DINOv3 ViT-L/16~\cite{dinov3_2025}
shared by all modalities; only the per-modality LoRA adapters
(rank~8), the fusion stack, the auxiliary decoders, and the
segmentation head are trained (46.8\,M trainable of 349.9\,M total,
13.4\%; the router adds ${\approx}0.27$\,M, \todo{exact full-model
parameter count from its training log}). Inputs are $1024^2$ crops.
We train for 200 epochs with AdamW (learning rate $6\times10^{-4}$,
weight decay 0.01), a warmup-poly schedule (power 0.9, 10-epoch
warmup), batch size 4 per GPU on 4 GPUs (effective batch 16 with
gradient accumulation), and bfloat16 mixed precision. The loss is
OHEM cross-entropy~\cite{ohem2016} plus the auxiliary
($\lambda_{\mathrm{aux}}{=}0.5$), calibration
($\lambda_{\mathrm{cal}}{=}0.1$), and router
($\lambda_{\mathrm{reg}}{=}0.01$, full model only) terms of
Sec.~\ref{sec:method}. Inference is single-scale without flipping or
any test-time augmentation. All experiments use a single seed (3407);
we therefore refrain from building claims on sub-1-point differences.

\textbf{Checkpoint-selection protocol.}
Both splits are evaluated every two epochs during training, which
creates a subtle risk of implicit test-set peeking when the reported
checkpoint is chosen post hoc. We adopt a strict protocol: all
headline numbers use the checkpoint with the best \emph{validation}
mIoU, and any test-selected (oracle) number is explicitly labeled as
such and excluded from state-of-the-art claims. Under this protocol
the no-router variant scores 68.19 val / 56.64 test (epoch~120); its
oracle test-best checkpoint reaches 57.60 (epoch~140) and is disclosed
only for completeness. One asymmetry remains: the no-router variant
was trained with a physics-based augmentation (PhysAug) that
competitors do not use, whereas the full model was trained without it;
\todo{PhysAug-free no-router re-run (T1) for a fully harmonized
ablation pair}.

\subsection{Comparison with State of the Art}
\label{sec:sota}

Table~\ref{tab:sota} compares ReliaDINO with published DELIVER
results. Following the protocol audit of Sec.~\ref{sec:related}, we
strictly separate the official \emph{test}-split cluster from
validation-only results and never mix the two; the ``extra
supervision'' column lists everything each method requires beyond
segmentation labels.

\begin{table}[t]
\centering
\caption{Comparison on DELIVER (four modalities: camera, LiDAR, depth,
event; 25 classes). ``Extra sup.''\ = supervision or data beyond
segmentation labels. Val and test columns are the official splits;
``--'' = not reported. $\dagger$: oracle test-selected checkpoint,
disclosed but excluded from claims (Sec.~\ref{sec:setup}).}
\label{tab:sota}
\setlength{\tabcolsep}{2.5pt}
\scriptsize
\resizebox{\columnwidth}{!}{%
\begin{tabular}{@{}llccc@{}}
\toprule
Method & Backbone & Extra sup. & Val & Test \\
\midrule
CMNeXt~\cite{cmnext2023}          & MiT-B2       & none                & 66.3  & 53.0 \\
GeminiFusion~\cite{geminifusion2024} & MiT-B2    & none                & 66.9  & 54.5 \\
MAGIC~\cite{magic2024}            & SegFormer-B2 & none                & 67.66 & --   \\
StitchFusion~\cite{stitchfusion2024} & MiT-B2    & none                & 68.18 & 53.4 \\
OmniSegmentor~\cite{omnisegmentor2025} & DFormer-L & mm.\ pretraining  & 68.0  & --   \\
CAFuser-CA$^2$~\cite{cafuser2025} & Swin-T       & cond.\ text (CLIP)  & 67.8  & 55.6 \\
CAFuser-CAA~\cite{cafuser2025}    & Swin-T       & cond.\ text (CLIP)  & 68.6  & 55.2 \\
DGFusion~\cite{dgfusion2026}      & Swin-T       & depth GT + cond.\ text & \todo{val n/a} & 56.7 \\
\midrule
ReliaDINO (no router)             & frozen DINOv3-L & none & 68.19 & 56.64\,/\,57.60$^\dagger$ \\
\textbf{ReliaDINO (full)}         & frozen DINOv3-L & none & \todo{67.74} & \todo{$\geq$57.14} \\
\midrule
\multicolumn{5}{@{}l@{}}{\todo{MUSES block (planned): DGFusion 79.5, CAFuser-CAA 78.5, ReliaDINO t.b.d.}}\\
\bottomrule
\end{tabular}%
}
\end{table}

On the test split, the full model reaches
\todo{57.14 (best checkpoint at the time of writing; training
completes on schedule and the final val-selected pair will replace
this number)}, exceeding the strongest published test result,
DGFusion at 56.7~\cite{dgfusion2026}, and the no-router variant's
oracle checkpoint reaches 57.60. Crucially, this is achieved with
segmentation labels \emph{only}: DGFusion additionally consumes
ground-truth log-depth supervision and CLIP-encoded condition
meta-text, and CAFuser requires condition annotations with
verbo-visual contrastive training. On the validation split, the
no-router variant (68.19) is within 0.41 of the best published result
(CAFuser-CAA, 68.6) while surpassing it by $+1.4$ on test. We note the
concurrent MM SAM-adapter~\cite{mmsamadapter2025}, which reports
strong DELIVER numbers (e.g., 57.35 RGB-D) under a \emph{two-modality}
protocol with a fine-tuned side network; our claims are scoped to the
four-modality, supervision-minimal setting. A residual caveat: an
eval-time resize-convention parity check against DGFusion is
outstanding, \todo{confirm eval-protocol parity (G0d) before quoting
head-to-head margins below ${\sim}2$ points as definitive}.

\subsection{Ablation Study}
\label{sec:ablation}

Table~\ref{tab:ablation} groups the ablations into training-time
variants (top) and inference-time module toggles applied to a fixed
checkpoint (bottom).

\begin{table}[t]
\centering
\caption{Ablations on DELIVER. Top: trained variants (val-selected
unless noted; $\dagger$ = oracle test-selected). Bottom:
inference-time toggles on the no-router test-best checkpoint,
full-resolution full test set. $^a$trained with PhysAug (see
Sec.~\ref{sec:setup}); $^b$run terminated by a system event at epoch
120/200; $^c$attention biases were active during this variant's
training but are inert at inference (toggle $\Delta{\approx}0$).}
\label{tab:ablation}
\setlength{\tabcolsep}{3pt}
\scriptsize
\resizebox{\columnwidth}{!}{%
\begin{tabular}{@{}lcc@{}}
\toprule
Trained variant & Val & Test \\
\midrule
ReliaDINO (full: router + gate + calibration) & \todo{67.74} & \todo{57.14} \\
\;\;no router, harmonized recipe$^b$ & 67.61 & 56.14 \\
\;\;no router, original recipe$^{a,c}$ & 68.19 & 56.64\,/\,57.60$^\dagger$ \\
\midrule
Inference toggle (base: test 56.64) & \multicolumn{2}{c}{$\Delta$ test mIoU} \\
\midrule
$-$ attention-logit reliability bias ($\lambda_1 B^{\mathrm{cal}}$) & \multicolumn{2}{c}{$|\Delta|\leq0.03$ (dead)} \\
$-$ consistency bias ($\lambda_2 B^{\mathrm{cons}}$) & \multicolumn{2}{c}{$|\Delta|\leq0.03$ (dead)} \\
$-$ competence gate + calibration & \multicolumn{2}{c}{$-0.26$ (56.64$\rightarrow$56.38)} \\
$-$ veto floor & \multicolumn{2}{c}{$+0.02$\,--\,$+0.23$ (below claim thresh.)} \\
\bottomrule
\end{tabular}%
}
\end{table}

\textbf{Router.}
The per-class reliability-anchored router is the difference between
the full model and the harmonized-recipe no-router variant
(same recipe, incomplete run): $+0.13$ val / $+1.00$ test at the
respective best checkpoints so far
(\todo{final router pair after training completes}). The router's
benefit concentrates on the test split and appears early: at matched
epochs it led the no-router model by $+2.0$ to $+2.4$ test mIoU
(epochs 30--46) and crossed the DGFusion test mark at epoch 58 versus
epoch 116 without the router. The comparison against the
original-recipe no-router variant is confounded by its PhysAug
training (footnote $a$), which is why we report both rows.

\textbf{Competence gate and calibration.}
Disabling the calibrated competence gate at inference costs $-0.26$
test mIoU on the full test set at full resolution. The veto floor
contributes only within our single-seed noise band and we do not claim
it. The train-time calibration loss is not separable by an inference
toggle; we claim only its diagnostic effect (balanced per-modality
reliability AUROC, Sec.~\ref{sec:analysis}), not an mIoU gain.

\textbf{Attention-logit reliability biasing: a negative result.}
Across five model generations spanning two backbone families (SAM2
memory attention and the present cross-modal attention), injecting
reliability as a pre-softmax additive bias on attention logits never
helped: on the current architecture the toggle changes test mIoU by at
most $0.03$ with fused-feature cosine similarity of $1.0$ (the bias is
functionally inert), and on the strongest SAM2-generation variant it
was a statistically significant net \emph{harm}
($p=4.5\times10^{-22}$, paired per-image test). We report this
placement finding as a contribution rather than hiding it: once
per-modality LoRA absorbs modality adaptation, reliability information
is useful where class decisions are formed (output gate and per-class
router), not inside the attention operator. Removing the bias also
restores the fused attention kernel at inference.

\textbf{Frozen-backbone choice.}
A controlled linear-probe on frozen backbones (identical head and
data) gives DINOv3 ViT-L 40.51 test mIoU on RGB versus 28.90 for the
SAM2 Hiera-B+ encoder ($+11.6$), and 35.48 versus 25.34 on depth
($+10.1$); on thin structures the gap is starker (TrafficLight IoU
32.2 vs.\ 8.27). Consistently, our four generations of SAM2-based
reliability designs gained only $+0.7$ val mIoU in total, whereas the
backbone swap alone contributed $+4.1$: fused-feature effective rank
rises from 1.26 to 10.18 and cross-modal CKA from 0.02--0.16 to
0.80--0.91, i.e., the frozen VFM removes the feature collapse that
earlier reliability mechanisms were implicitly compensating for.

\textbf{Per-modality LoRA.}
Disabling a modality's LoRA adapters (frozen-backbone features
instead) drops that modality's auxiliary pixel accuracy in all five
conditions: by $+0.157$ to $+0.202$ for LiDAR and $+0.110$ to $+0.153$
for depth (the modalities farthest from the RGB pretraining
distribution), and by $+0.032$ to $+0.064$ (RGB) and $+0.020$ to
$+0.050$ (event). All 48 adapter pairs (24 blocks $\times$ 2
projections) remain active for every modality (0/48 dead), confirming
that the single frozen encoder is genuinely multiplexed rather than
defaulting to its RGB prior.

\subsection{Per-Class and Per-Condition Analysis}
\label{sec:analysis}

\textbf{Per-condition test audit.}
On the no-router test-best checkpoint, per-condition test mIoU is
55.64 (cloud), 56.32 (fog), 54.56 (night), 56.31 (rain), and 55.41
(sun) --- a spread of only 1.76 points. The large val--test gap of
DELIVER is therefore \emph{not} an adverse-condition effect but a
per-class transfer effect. To our knowledge this is the first
published per-class, per-condition audit of the DELIVER \emph{test}
split; competitors publish only aggregate test numbers, so
Table~\ref{tab:sota} has no external per-condition column.

\textbf{Benchmark ceiling.}
Three classes are effectively dead for \emph{all} methods we could
audit: Other (max IoU 5.6 across conditions), Wall (8.7), and Bridge
(0.5) score near zero for our model, for the frozen DINOv3 linear
probe, and for the official CMNeXt checkpoint
(\todo{exact official CMNeXt per-class test numbers from the M0-a
audit artifact}). We flag these as a benchmark ceiling: no fusion
mechanism can be credited or blamed for them. Conversely, classes that
were dead in our SAM2-generation models are recovered by the present
design: Water $0.0\rightarrow5.4$, TrafficLight $12.4\rightarrow36.6$,
Static $29.4\rightarrow39.1$ (five-condition means).

\textbf{Reliability quality.}
The calibration objective yields \emph{balanced} per-modality
reliability: at night, the AUROC of the calibrated confidence as a
correctness predictor is 0.85 / 0.78 / 0.87 / 0.70 for RGB / depth /
event / LiDAR. Earlier generations either collapsed on non-RGB
modalities (e.g., 0.84 / 0.63 / 0.26 / 0.38) or repaired non-RGB
geometry at the cost of RGB (0.43 / 0.92 / 0.55 / 0.97); ours is the
only generation with all four modalities above 0.70 without repair.

\textbf{Modality contributions.}
Drop-modality deltas (mIoU loss when a modality is removed at
inference) show depth as the dominant complement ($+9.5$ to $+11.4$
across conditions) and RGB rising from $+2.8$ (cloud) to $+8.0$
(night); LiDAR adds ${\leq}1.2$, and event contributes ${\approx}0$ in
all five conditions --- a known, lineage-wide limitation that fusion
alone does not fix. Qualitative results under adverse conditions are
shown in Fig.~\ref{fig:qual}, and the per-class audit, AUROC balance,
and gating behavior are visualized in Fig.~\ref{fig:analysis}.

\subsection{Fairness Protocol}
\label{sec:fairness}

For a clean reading of Table~\ref{tab:sota} we disclose all known
asymmetries. \emph{Ours:} a larger frozen backbone (349.9\,M total
parameters), though our 46.8\,M trainable parameters are fewer than
the ${\sim}58.7$\,M \emph{total} of CMNeXt-B2 and comparable to
trainable Swin-T pipelines (${\sim}79$\,M); the no-router variant used
PhysAug (Sec.~\ref{sec:setup}, \todo{harmonized re-run}); the full
model does not. \emph{Theirs:} DGFusion trains with ground-truth
log-depth supervision and CLIP condition meta-text, and CAFuser with
condition annotations; both use roughly a $2\times$ longer effective
training schedule (ours: 200 epochs). \emph{Both directions:} no
test-time augmentation on our side, matched input modalities, and
val-only checkpoint selection for every headline claim.


% ---------------- Fig. 4: qualitative grid (two-column, placeholder) -------
\begin{figure*}[t]
  % ============================================================
% fig:qual — qualitative comparison grid (PLACEHOLDER, assets TODO)
% BODY-ONLY: \input inside a two-column {figure*} environment.
% Planned grid: rows = conditions (night / fog / rain / corruption),
% cols = RGB input / GT / baseline (CMNeXt or CAFuser) / ReliaDINO.
% Asset sources (see figures/figure_plan.md, Fig. 4):
%  - ours: tools/eval_per_domain.py masks; existing no-router masks at
%    /drone_nas/drone/analysis_logs/P34_eval_20260713/per_domain/masks_best_<cond>/
%  - baseline predictions: \todo (no local baseline checkpoint yet)
% ============================================================
\centering
\setlength{\fboxsep}{0pt}%
\providecommand{\qcell}{}\renewcommand{\qcell}[1]{\framebox{\parbox[c][2.2cm][c]{0.225\textwidth}{\centering\scriptsize #1}}}%
\begin{tabular}{@{}c@{\;}c@{\,}c@{\,}c@{\,}c@{}}
 & {\scriptsize RGB input} & {\scriptsize Ground truth} &
   {\scriptsize \todo{CMNeXt or CAFuser}} & {\scriptsize ReliaDINO (ours)} \\
\rotatebox{90}{\scriptsize night} &
 \qcell{\todo{night RGB}} & \qcell{\todo{night GT}} &
 \qcell{\todo{night baseline}} & \qcell{\todo{night ours}} \\
\rotatebox{90}{\scriptsize fog} &
 \qcell{\todo{fog RGB}} & \qcell{\todo{fog GT}} &
 \qcell{\todo{fog baseline}} & \qcell{\todo{fog ours}} \\
\rotatebox{90}{\scriptsize rain} &
 \qcell{\todo{rain RGB}} & \qcell{\todo{rain GT}} &
 \qcell{\todo{rain baseline}} & \qcell{\todo{rain ours}} \\
\rotatebox{90}{\scriptsize corrupt.} &
 \qcell{\todo{corruption RGB}} & \qcell{\todo{corruption GT}} &
 \qcell{\todo{corruption baseline}} & \qcell{\todo{corruption ours}} \\
\end{tabular}
\caption{Qualitative results on the DELIVER test split under adverse
conditions (rows: night, fog, rain, and a sensor-corruption case). Columns
show the RGB input, ground truth, the prediction of
\todo{CMNeXt/CAFuser baseline with checkpoint provenance}, and ReliaDINO.
ReliaDINO preserves thin structures (poles, road lines, traffic lights) and
remains stable where the RGB signal degrades, consistent with the
reliability-driven gate and per-class router.
\todo{Insert final-checkpoint predictions; add class-color legend strip.}}
\label{fig:qual}

\end{figure*}

% ---------------- Fig. 5: analysis panels (two-column, placeholder) --------
\begin{figure*}[t]
  % ============================================================
% fig:analysis — 2-panel analysis figure (PLACEHOLDER, plots TODO)
% BODY-ONLY: \input inside a two-column {figure*} environment.
% Panel (a): per-class DELIVER test IoU audit (25 classes, 5-condition
%   mean), dead classes (Other/Wall/Bridge) highlighted.
%   Data: /drone_nas/drone/analysis_logs/P34_eval_20260713/per_domain_analysis.md
%   (tools/analyze_per_domain.py output).
% Panel (b): reliability AUROC per modality at night + drop-modality
%   ΔmIoU per condition.
%   Data: /drone_nas/drone/analysis_logs/P34_eval_20260713/module_diag.json
%   (tools/module_diagnostics.py / tools/eval_reliability_auroc.py).
% Export script TODO: _paper_submission/figures/scripts/ (matplotlib → PDF).
% ============================================================
\centering
\setlength{\fboxsep}{0pt}%
\framebox{\parbox[c][5.2cm][c]{0.48\textwidth}{\centering\scriptsize
  (a) per-class DELIVER \emph{test} IoU audit\\[2pt]
  \todo{horizontal bar chart: 25 classes sorted by 5-condition-mean IoU;\\
  Other/Wall/Bridge hatched red with ``dead for official CMNeXt too'' bracket}}}%
\hfill
\framebox{\parbox[c][5.2cm][c]{0.48\textwidth}{\centering\scriptsize
  (b) reliability AUROC (night) + drop-modality $\Delta$mIoU\\[2pt]
  \todo{top: 4 bars, AUROC .85/.78/.87/.70 (RGB/depth/event/LiDAR),\\
  chance line at 0.5; bottom: grouped bars per condition,\\
  depth $\approx{+}9.5$--$11.4$, event $\approx 0$ annotated}}}
\caption{Analysis on the DELIVER test split. (a)~Per-class audit
(five-condition mean IoU): the residual gap to the validation split is
concentrated in a small set of classes; \emph{Other}, \emph{Wall}, and
\emph{Bridge} remain near zero for our model, for a linear probe on the
frozen backbone, and for the official CMNeXt model
\todo{per-class CMNeXt test numbers (M0-a audit)}, indicating a benchmark
ceiling rather than a fusion failure. (b)~Left: per-modality reliability
AUROC at night (0.85/0.78/0.87/0.70 for RGB/depth/event/LiDAR) --- the
calibrated signals rank correctness well for all four modalities. Right:
drop-modality $\Delta$mIoU per condition; depth is the dominant complement
(${\approx}{+}9.5$--$11.4$), whereas event contributes ${\approx}0$ in all
five conditions. Diagnostics computed on the no-router variant's audited
checkpoint; \todo{regenerate for the final full model if time permits}.}
\label{fig:analysis}

\end{figure*}

% ============================================================
% Section 5 — Conclusion (body-only; \input from root.tex)
% Numbers: fact_experiments.md (P36 final = TODO until run completes)
% ============================================================

\section{Conclusion}
\label{sec:conclusion}

We presented ReliaDINO, a multimodal semantic segmentation model that
multiplexes a \emph{single frozen} DINOv3 ViT-L/16 across four sensing
modalities through per-modality LoRA adapters, exchanges information with two
cross-modal attention layers, and fuses the modalities with a calibrated
competence gate and a per-class reliability-anchored router. Trained with
segmentation labels only---no depth supervision, no condition
annotations---ReliaDINO exceeds the previous best test-set result on
DELIVER~\cite{cmnext2023} held by the more heavily supervised
DGFusion~\cite{dgfusion2026}
(\todo{final full-model test mIoU after training completes; best observed
57.14 vs.\ 56.71, and 57.60 for the no-router variant under test-side
selection / 56.64 under val-only selection}). Beyond the benchmark numbers,
our controlled five-generation audit yields a placement lesson for
reliability-aware fusion: self-derived reliability signals are actionable
\emph{where class decisions are formed}---at the output gate and the
per-class router---whereas injecting the same signals as pre-softmax
attention-logit biases becomes a no-op, or a measurable harm, once
per-modality low-rank adaptation of a strong frozen backbone has already
aligned the modality representations.

Limitations remain. Our evidence is so far confined to a single benchmark
and a single seed, so we refrain from claims below one mIoU point; and the
frozen-backbone recipe inherits a per-class ceiling that official baselines
share (Wall, Bridge, and Other stay near zero for our model, for a linear
probe of the backbone, and for the official CMNeXt
model~\cite{cmnext2023}\todo{confirm official per-class test source (M0-a)
before final wording}). Future work extends the recipe to
MUSES~\cite{muses2024} and to aquatic scenes (MULTIAQUA), adds multi-seed
verification, and transfers reliability-routed fusion from dense
segmentation to multimodal detection.

\section*{Acknowledgments}
\todo{Funding sources and acknowledgments.}


% \addtolength{\textheight}{-2cm} % balance last-page columns before submission

%%%%%%%%%%%%%%%%%%%%%%%%%%%%%%%%%%%%%%%%%%%%%%%%%%%%%%%%%%%%%%%%%%%%%%%%%%%%%%
\bibliographystyle{IEEEtran}
\bibliography{references}

\end{document}
